\section{System Overview}
\label{sec:overview}

Figure~\ref{fig:pipeline} shows the four-stage pipeline---\emph{generate,
extract, verify, source-or-refuse}---and Algorithm~\ref{alg:pipeline} states it
precisely. A citizen question $q$ is sent to a sovereign LLM
(\S\ref{sec:sovereign}), which returns a draft answer $a$. The draft is parsed
for legal citations of the canonical French form \emph{``article
$\langle$num$\rangle$ du $\langle$code$\rangle$''} by a normalization-tolerant
extractor (\S\ref{sec:verify}). Each extracted citation $c_i$ is verified
\emph{independently} against the legal database, yielding a verdict together
with, on success, a canonical \texttt{LEGIARTI} identifier, a clickable
Légifrance URL, and a short excerpt.

\begin{algorithm}[t]
\caption{Fail-closed answering. A single unverifiable citation---or any
exception raised while checking---yields a refusal.}
\label{alg:pipeline}
\begin{algorithmic}[1]
\Function{Answer}{$q$}
  \State $a \gets \textsc{Llm}(q)$ \Comment{untrusted draft}
  \State $C \gets \textsc{ExtractCitations}(a)$
  \If{$C = \emptyset$}
    \State \Return \textsc{Refuse}
  \EndIf
  \State $S \gets \emptyset$
  \For{$c \in C$}
    \State $r \gets \textsc{Verify}(c)$ \Comment{state-aware DB lookup}
    \If{$\neg\, r.\mathit{exists}$}
      \State \Return \textsc{Refuse} \Comment{fail-closed}
    \EndIf
    \State $S \gets S \cup \{\langle c,\, r.\mathit{url},\, r.\mathit{excerpt}\rangle\}$
  \EndFor
  \State \Return \textsc{Sourced}(a, S)
\EndFunction
\end{algorithmic}
\end{algorithm}

\paragraph{The release predicate.}
Let $C(a) = \{c_1, \dots, c_n\}$ be the citations extracted from draft $a$. The
answer is released to the user \emph{if and only if}
\[
  C(a) \neq \emptyset \;\wedge\; \forall\, c_i \in C(a):\;
  \mathrm{verify}(c_i) = \textsc{ok}.
\]
Every other outcome---no citation at all, one unresolved citation, or any
exception while checking---maps to an explicit refusal. The disjunction is
deliberately lopsided: the release branch demands \emph{universal} success,
whereas a \emph{single} failure anywhere triggers refusal. This is what we mean
by \emph{fail-closed}, and it is the whole point: the guarantee does not depend
on the model behaving well, only on the verifier being sound about what it can
confirm.

\paragraph{Why independent, per-citation checks.}
Verifying citations one at a time---rather than scoring the answer as a
whole---keeps the trusted computing base small and its verdicts explainable:
each released citation is accompanied by the exact identifier and URL that
justify it, so a user (or an auditor) can click through to the primary source.
Independence also makes the layer trivially parallelizable and lets a partial
failure localize to the offending citation instead of tainting the whole
answer.

\begin{figure}[t]
\centering
\begin{tikzpicture}[
  node distance=6mm and 8mm,
  box/.style={draw, rounded corners, align=center, font=\footnotesize,
              minimum height=7mm, inner sep=3pt},
  db/.style={draw, cylinder, shape border rotate=90, aspect=0.25,
             align=center, font=\footnotesize, inner sep=3pt},
  dec/.style={draw, diamond, aspect=2, align=center, font=\scriptsize, inner sep=1pt},
  ok/.style={draw=lrblue, thick, rounded corners, align=center, font=\footnotesize, inner sep=3pt},
  no/.style={draw=lrred, thick, rounded corners, align=center, font=\footnotesize, inner sep=3pt},
  a/.style={-{Latex}, thick}]
  \node[box] (q) {Question};
  \node[box, right=of q] (g) {LLM\\(sovereign)};
  \node[box, below=of g] (e) {Extract\\citations};
  \node[dec, below=of e] (v) {verify?};
  \node[db, left=of v, xshift=-2mm] (db) {Canutes\\Légifrance};
  \node[ok, right=of v, xshift=6mm] (r) {Sourced answer\\+ Légifrance link};
  \node[no, below=of v, yshift=-1mm] (x) {Refusal\\``no source,\\no answer''};
  \draw[a] (q) -- (g);
  \draw[a] (g) -- (e);
  \draw[a] (e) -- (v);
  \draw[a] (v) -- (db) node[midway, above, font=\scriptsize] {num+code};
  \draw[a] (db) -- (v) node[midway, below, font=\scriptsize] {VIGUEUR?};
  \draw[a] (v) -- (r) node[midway, above, font=\scriptsize] {all exist};
  \draw[a] (v) -- (x) node[midway, right, font=\scriptsize] {any missing};
\end{tikzpicture}
\caption{The fail-closed pipeline: \emph{generate, extract, verify,
source-or-refuse}. The trust guarantee is enforced at the verification
diamond, independently of the model.}
\label{fig:pipeline}
\end{figure}

\subsection{Trusted computing base}
The components that must be correct for the guarantee to hold are few and
explicit: the citation extractor, the verifier, and the legal database. The
model, the network, the serving hardware, and even the prompt lie
\emph{outside} the trusted base---any may fail arbitrarily, and the worst
observable outcome is a refusal, never a fabricated citation presented as
verified. Shrinking the trusted base to a database lookup plus a few hundred
lines of verification code is what makes the guarantee auditable rather than
aspirational.

\subsection{Where the latency goes}
End-to-end latency is dominated by \emph{generation}, not verification: the
verifier issues a handful of bounded, index-friendly queries
(\S\ref{sec:verify}) whose cost is negligible beside a multi-hundred-token
decode. Per-citation checking is therefore affordable online, and the
end-to-end latency we report (\S\ref{sec:eval}) tracks the model rather than
the database---verification buys the guarantee essentially for free relative to
the cost already paid to generate the draft.
